\documentclass[11pt,a4paper]{article}
\usepackage[UTF8]{ctex}
\usepackage{geometry}
\geometry{left=2.5cm,right=2.5cm,top=2.5cm,bottom=2.5cm}
\usepackage{amsmath,amssymb,amsthm}
\usepackage{hyperref}
\usepackage{booktabs}
\usepackage{enumitem}
\usepackage{xltabular}
\hypersetup{colorlinks=true,linkcolor=blue,urlcolor=blue}

\newtheorem{definition}{定义}[section]
\newtheorem{proposition}{命题}[section]
\newtheorem{remark}{注记}[section]

\title{攻防博弈论统一框架（修正版）\\
——从探测博弈到AI战场转移的均衡分析}
\author{李龙胜}
\date{\today}

\begin{document}

\maketitle

\begin{center}
\textit{
探测与反探测是最基础的信号博弈。\\
逼移是连续时间上的动态消耗战——\\
当前从开环起步，真正的战场在闭环。\\
多攻击方协同是合作博弈的核心分配问题。\\
AI战场转移将弹药质量变为新的均衡维度。\\
四个层面，同一套博弈模板——\\
防御方私有信息，攻击方不完全观测，\\
双方贪心择优，边际均衡统治一切。\\
本文是攻防博弈论的统一场论：\\
它把分散在四卷和扩展模块中的博弈工具全部显式化，\\
揭示所有攻防分析共享同一个博弈结构。
}
\end{center}

\vspace{5mm}

\begin{abstract}
攻防双方的对抗过程本质上是信息不对称下的动态零和博弈。
本文在前序四卷体系和扩展模块的基础上，
将分散于探测、逼移、多攻击方协同和AI战场转移四个层面的博弈结论
统一为一个完整的深化框架。
基础探测博弈被精确化为信号博弈，
逼移策略被建模为连续时间最优控制问题，
多攻击方协同被纳入合作博弈的Shapley值分配，
AI战场转移揭示了弹药质量取代弹药数量成为新均衡维度的深层机制。
在外部审查基础上，补充了连续类型空间下精炼贝叶斯均衡的验证需求、
逼移动态博弈的开环与闭环框架差异、
以及AI弹药质量均衡的时序假设讨论。
四个层面共享同一博弈模板，
贪心择优定理在不完全信息动态博弈中的推广
是整个体系的核心数学基础。
\end{abstract}

\tableofcontents

\section{引言：博弈论的统一需求}

在已有的安全运营理论体系中，
攻防博弈的结论分散于多个文档——
第2卷的探测与反探测、第3卷的逼移策略与均衡分析、
问题六的攻击方联盟信任锚定与瓦解、
以及AI化攻防专题卷和批判文档中的战场转移分析。
这些文档各自独立，但共享同一个博弈结构。

本文将这些分散的博弈结论整合为一个统一的深化框架。
四个层面的博弈——探测与反探测、逼移资源消耗战、
多攻击方协同与竞争、AI战场转移——
被严格形式化，并揭示它们在数学结构上的同构性。
本文在理论体系中承担一个独特的角色：
它把四卷体系和所有扩展模块中“隐含使用”的博弈论工具全部显式化，
让读者看到所有攻防分析共享同一个博弈模板——
不完全信息动态博弈下的贪心均衡。

\section{基础博弈模板：探测与反探测的信号博弈}

这是整个攻防博弈的原子单元。

\subsection{参与者与私有信息}

防御方持有私有类型 \(\theta\)（秘密参数），
攻击方持有先验信念 \(\pi(\theta)\)。
防御方知道 \(\theta\)，攻击方只知道 \(\pi(\theta)\)。
这是典型的贝叶斯博弈设定。

在安全运营实例中，防御方的 \(\theta\) 是一个高维连续空间——
每个告警类别的 \(A_i^{\min}\)、\(\alpha_i\)、扰动幅度 \(\delta\) 等参数。
攻击方的信念 \(\pi(\theta)\) 是这个高维空间上的概率分布。

\subsection{时序}

博弈分为三个阶段：
\begin{enumerate}[label=(\arabic*)]
    \item 防御方选择分析率向量 \(\mathbf{r} = (r_1,\dots,r_K)\)
          和伪造策略 \((q_{01}, q_{10})\)。
    \item 攻击方选择探测序列，对每类告警发起 \(n_i\) 次探测，
          观测响应 \(y_i \in \{0,1\}\)。
    \item 攻击方基于观测更新信念 \(\pi(\theta \mid \mathbf{y})\)，
          并选择攻击方向 \(i_A\) 和攻击量。
\end{enumerate}

\subsection{收益函数}

防御方最小化漏报损失加分析成本加伪造成本，
攻击方最大化防御方的漏报损失减去探测成本。
双方零和。

\subsection{均衡特征}

这是一个信号博弈。
防御方的策略是信号（分析率和伪造参数的选择），
攻击方的策略是响应（探测和攻击方向的选择）。
贝叶斯-纳什均衡中：
防御方的策略使得攻击方无法通过改变探测策略来提高期望收益；
攻击方的策略使得防御方无法通过改变分析率来降低期望损失。

\subsection{伪造在均衡中的作用}

防御方通过伪造 \((q_{01}, q_{10})\) 使攻击方观测到的
\(P_{\text{obs}}(y=1)\) 偏离真实 \(r_i\)。
攻击方知道伪造的存在，但不知道具体参数，
因此只能做稳健估计。
均衡时，防御方的伪造强度正好抵消攻击方的信息增益，
使得攻击方的最优估计精度被锁定在某个下界——
这就是信息论下界的博弈论版本。

最优伪造强度不是越大越好，
而是正好使攻击方的边际信息增益等于防御方的边际伪造成本。
这与够用就行原则完全一致。

\textbf{精炼贝叶斯均衡的验证需求}：
当前模型描述了信号博弈的均衡结构，
但精炼贝叶斯均衡在连续类型空间和高维信号空间中的
存在性和唯一性，需要进一步的数学验证。
防御方的 \(\theta\) 是高维连续空间，
攻击方的信念 \(\pi(\theta)\) 是该空间上的概率分布，
在这种设定下PBE的存在性不自动成立。
当前均衡特征的定性描述提供了直观理解，
严格数学证明是后续工作的重要方向。

\section{逼移策略的动态均衡——连续时间最优控制}

逼移不是一次性的，而是在时间上演化的。
将其建模为一个连续时间动态博弈。

\subsection{状态演化}

设防御方分析能力消耗为 \(C(t)\)，总容量为 \(C_{\max}\)。
攻击方在逼移方向 \(i_b\) 施加压力 \(p(t)\)，
在真实攻击方向 \(i_a\) 施加攻击强度 \(a(t)\)。

防御方分析能力的演化：
\[
\frac{dC}{dt} = p(t) \cdot \kappa_b + a(t) \cdot \kappa_a,
\]
其中 \(\kappa_b, \kappa_a\) 分别是逼移方向和真实攻击方向的单次分析成本。
当 \(C(t) \to C_{\max}\) 时，防御方进入容量极限，
分析响应延迟骤增，封禁策略从精确封禁相变为批量封禁。

\subsection{收益与最优控制}

攻击方在时间窗口 \([0, T]\) 内选择 \(p(t)\) 和 \(a(t)\)，
最大化终端时刻 \(T\) 真实攻击的成功概率，
减去逼移成本的累积：
\[
\max_{p(t), a(t)} \left[ a(T) \cdot A_{\text{target}} - \int_0^T c_b \cdot p(t)\,dt \right].
\]
防御方选择各路分析率 \(r_i(t)\)，受容量约束：
\[
\min_{\mathbf{r}(t)} \left[ (1 - r_a(T)) \cdot a(T) \cdot A_{\text{target}} \right],
\quad \text{s.t.} \quad \sum_i r_i(t) \cdot N_i \cdot \kappa_i \le C_{\max} - C(t).
\]

\subsection{逼移最优时机的边际条件}

这是一个微分博弈。
它的鞍点均衡给出逼移最优时机：
当防御方能力饱和度 \(s(t) = C(t)/C_{\max}\)
达到某个临界值 \(s^*\) 时，
逼移的边际效果突然放大——
防御方分析能力曲线的凸性在 \(s^*\) 处急剧增加，
维持正常分析的边际成本骤增，贪心稳态失效。
攻击方在此时集中释放逼移压力，效果最大化。

这一位相变的结构与引力场生成论推导中
贪心跳变点的结构完全同构——
当累积成本接近容量极限时，
贪心失效是数学必然，不是偶然。

\textbf{开环框架与闭环框架的差异}：
当前模型采用了开环最优控制框架——
攻击方在时间零点规划好整个窗口内的 \(p(t)\) 和 \(a(t)\)。
在真实攻防中，双方持续观测对方行为并实时调整策略：
攻击方观察到防御方响应延迟后调整逼移强度，
防御方观察到攻击方逼移模式后调整分析率分配。
应采用闭环反馈纳什均衡框架。
闭环解和开环解在微分博弈中通常存在质的差异——
逼移最优时机的临界饱和度 \(s^*\) 可能随
防御方的反馈策略而动态偏移。
当前开环解是闭环均衡的一阶近似，
更精确的闭环分析是后续研究的重要方向。

\section{多攻击方协同的合作博弈解——信任锚定的边际成本}

当攻击方不止一个时，博弈从非合作升级为合作博弈。

\subsection{联盟价值函数}

设攻击方联盟 \(\mathcal{M}\) 协同攻击带来的总收益为
\(v(\mathcal{M})\)——防御方漏报损失的增幅。
这个增幅源于信息共享、分工协同和轮班消耗
对防御方分析能力的联合压制。

\subsection{Shapley值与核心解}

联盟总收益必须在成员之间分配，
每个成员的Shapley值等于该成员对所有可能子联盟的边际贡献的平均。
若分配方案满足个体理性和联盟理性，且不存在任何子联盟
通过退出联盟可以获得更高收益，则联盟稳定。
这个稳定解就是合作博弈的核心。

\subsection{信任锚定与违约成本}

去中心化自治协同中的抵押品和智能合约，
对应合作博弈中的可信承诺机制。
抵押品被罚没的风险等价于退出联盟的违约成本。
当违约成本大于退出联盟的收益增额时，
联盟稳定在核心解中。
当违约成本不足以约束背叛时，
联盟面临囚徒困境。

\subsection{防御方的瓦解策略}

防御方通过渗透进联盟制造虚假信号——
伪造某些成员的贡献数据，
使联盟内部的Shapley值分配偏离真实贡献，
引发分配纠纷。
这恰好是分而治之策略在合作博弈框架下的精确表述。

\textbf{联盟价值函数的精确参数化需求}：
Shapley值的计算依赖于联盟价值函数 \(v(\mathcal{M})\) 的精确定义。
该函数依赖于各攻击方的资源预算、私有信息、
防御方对联盟攻击的响应策略、以及联盟内部的信息共享效率。
在真实场景中这些参数极难获取。
当前版本给出的Shapley值框架是概念性的——
它在定性层面成立，说明联盟收益存在且可分配。
精确数值解需要攻击方资源数据和防御方响应策略的完整描述，
待实证数据标定。

\section{AI战场转移——弹药质量成为均衡维度}

当攻防双方都引入大模型时，
博弈的战场从网络流量层精确转移到模型参数空间。

\subsection{战场转移的精确坐标}

五对战场坐标的转移对应表（见AI化攻防批判文档）
在此被纳入博弈论框架。
每一对战场坐标在新坐标系中保持了完全相同的数学结构。
转移的不是博弈规则，而是博弈空间的坐标轴。

\subsection{弹药质量的纳什均衡}

AI时代的根本转变在均衡维度中找到了精确的表达。

设攻击方生成对抗样本的质量为 \(q_A\)——逃逸检测的概率，
防御方模型训练的质量为 \(q_D\)——检测率。
两者的成本函数分别为 \(C_A(q_A)\) 和 \(C_D(q_D)\)，
均严格凸。

攻击方的收益：\(q_A \cdot A_{\text{target}} - C_A(q_A)\)。
防御方的收益：\(-(1 - q_D(q_A)) \cdot A_{\text{target}} - C_D(q_D)\)，
其中 \(q_D(q_A)\) 是防御方针对攻击方质量 \(q_A\)
进行对抗训练后的最终检测率。

双方在质量维度上形成纳什均衡：
攻击方选择 \(q_A^*\) 使边际收益等于边际生成成本，
防御方选择 \(q_D^*\) 使边际检测提升等于边际训练成本。
弹药质量取代弹药数量成为新的均衡维度。

\subsection{AI时代的深刻悖论}

AI引入本应降低防御方的边际分析成本，
但质量竞赛的加剧可能反而增加防御方的总成本——
因为维持一个能抵抗高质量对抗样本的模型
需要持续、高成本的对抗训练。
防御方在AI时代的竞争可能比传统时代更激烈。

但够用就行原则在此同样有效：
防御方在质量维度上同样要寻求“刚好够用”的模型精度，
而非追求完美。
最优模型精度 \(q_D^*\) 恰好是
边际检测收益等于边际训练成本的均衡点。

\textbf{弹药质量均衡的时序假设}：
当前模型的静态纳什均衡假设双方同时行动。
在真实AI攻防中，时序可能是交替的——
防御方先训练模型（选择初始 \(q_D\)），
攻击方观测检测率后生成对抗样本（选择 \(q_A\)），
防御方再观测对抗样本并进行对抗训练（更新 \(q_D\)）。
这是一个交替行动的动态博弈。
如果防御方先动，可能通过先发优势锁定更高的 \(q_D^*\)；
如果攻击方能够完全观测防御方模型精度后再行动，
则拥有后发优势。
在持续对抗训练中，双方交替行动，动态收敛到稳态均衡。
当前模型的静态纳什均衡是这一动态稳态的一阶近似，
更多样的时序设定下的均衡比较是后续研究的重要方向。

\section{统一博弈模板——正在浮现的核心定理}

以上四个层面共享同一个博弈模板：
防御方拥有私有信息 \(\theta\)，
攻击方拥有不完全观测 \(y\)，
双方各自拥有有限资源预算，
双方在每一步选择当前最优操作，
信息不对称是防御方的天然优势——
但这一优势依赖于攻击方探测成本不能太低。

这个模板的数学本质是贪心择优定理
在不完全信息动态博弈中的推广——
当成本函数满足凸性条件且进展量单调时，
每步贪心选择等价于全局最优响应。
在所有四个层面上，
防御方和攻击方的最优策略都可以由贪心择优统一裁定。

这一推广如果被严格证明——
贪心择优在不完全信息动态博弈中等价于贝叶斯-纳什均衡——
将使整个四卷体系在博弈论维度上获得I级的严格性。
当前该定理的推广骨架清晰，
严格数学证明是未来最重要的理论攻坚方向之一。

\section{诚实标注}

\begin{center}
\begin{xltabular}{\textwidth}{c|X|c}
\toprule
\textbf{命题} & \textbf{状态} & \textbf{层级} \\
\midrule
探测与反探测信号博弈 &
形式化完整，均衡定性结论明确，
最优伪造强度的边际条件已给出 &
II（精炼贝叶斯均衡在高维连续类型空间中的存在性和唯一性需进一步数学验证） \\
逼移动态均衡——连续时间最优控制 &
微分博弈框架已建立，
贪心跳变点的位相变条件已推导 &
II（开环框架是闭环均衡的一阶近似，精确临界值 \(s^*\) 需闭环分析） \\
多攻击方协同的合作博弈解 &
Shapley值与核心解框架已引入，
信任锚定与违约成本的对应关系已建立 &
II（联盟价值函数 \(v(\mathcal{M})\) 的精确参数化需实证数据） \\
AI战场转移——弹药质量均衡 &
质量维度的纳什均衡已推导，
AI时代的深层悖论已论证 &
II（静态均衡是动态交替行动稳态的一阶近似，时序假设的影响需进一步分析） \\
统一博弈模板——贪心择优定理的博弈推广 &
四个层面的同构性已严格论证，
定理推广骨架清晰 &
II（严格推广证明需进一步的数学工作） \\
\bottomrule
\end{xltabular}
\end{center}

\section{结语}

攻防博弈论统一框架的深化修正版本已完成。
探测与反探测是最基础的信号博弈，
逼移是连续时间上的动态消耗战——
当前从开环起步，真正的战场在闭环。
多攻击方协同是合作博弈中的核心分配问题，
AI战场转移将弹药质量变为新的均衡维度。

四个层面共享同一个博弈模板：
防御方私有信息，攻击方不完全观测，
双方贪心择优，边际均衡统治一切。
贪心择优定理在不完全信息动态博弈中的推广，
是贯穿整个体系的核心数学基础——
一旦严格证明，它将使四卷体系在博弈论维度上获得I级的严格性。

\vspace{5mm}
\noindent\textbf{文档信息}：
\begin{itemize}
    \item 作者：李龙胜
    \item 版本：修正版（整合外部审查反馈）
    \item 许可：CC BY 4.0
    \item 前置依赖：四卷体系、问题六修正版、
          AI化攻防专题卷、AI化攻防批判修正版、
          引力场与黑洞的生成论推导
\end{itemize}

\end{document}